%% kapa: exam=1 round=29 year=2018 period=1 form=A subject=real_estate subject_ko=부동산학원론 seq=14 number=94
%% source: Q-Net 감정평가사 기출문제 (cst00309) artlSeq=5202808 file=extracted/1008/5202808/2018년 제29회 감정평가사 1차시험 기출문제/2018년 제29회 감정평가사 1차시험 기출문제/1교시 A형.pdf
%% conversion: kordoc 4.24.51 -> tools/gen1.py (auto); review=proofread-sonnet; flags=none
\begin{question}{94}
포트폴리오 이론에 따른 부동산투자의 포트폴리오 분석에 관한 설명으로 옳지 않은 것은?
\begin{choices}
  \item 체계적 위험은 분산투자를 통해서도 회피할 수 없다.
  \item 위험과 수익은 상충관계에 있으므로 효율적 투자선은 우하향하는 곡선이다.
  \item 투자자의 무차별곡선과 효율적 투자선의 접점에서 최적의 포트폴리오가 선택된다.
  \item 비체계적 위험은 개별적인 부동산의 특성으로 야기되며 분산투자 등으로 회피할 수 있다.
  \item 포트폴리오 구성자산의 수익률 간 상관계수(ρ)가 ‘-1’인 경우는 상관계수(ρ)가 ‘1’인 경우에 비해서 위험 회피효과가 더 크다.
\end{choices}
\end{question}
\answer{2}
